\section{\uppercase{Related Work}}
\label{sec:related}
\noindent Natural-language visualization already separates interpretation from graphical specification. NL4DV maps data and queries to attributes, analytical tasks, and Vega-Lite specifications \cite{nl4dv}. NL4DV-LLM adds prompting with ambiguity handling and conversation \cite{sah2024}; its error analysis includes missing tasks and misleading attribute bindings. Neither structured representation nor the existence of omissions is a novelty claim here.\draftassist

LLM-based visualization systems vary in what they delegate. LIDA combines data summarization, goal exploration, generation and evaluation of visualization code, and infographic generation \cite{lida}. Data Formulator~2 combines graphical and natural-language inputs, generates data transformations, and preserves branches of an iterative authoring history \cite{dataformulator2}. These systems establish broader authoring and interaction capabilities than the fixed analyses used here. Our model neither writes analysis programs nor designs arbitrary layouts. The question is reliability of answer selection when computation and rendering are controlled. The saved-reference selector only reads existing results.

InsightBench evaluates multi-step analytical agents that formulate questions, analyze answers, and summarize insights using datasets with planted patterns \cite{insightbench}. Our comparison addresses explicit longitudinal questions on unmodified observations and tracks which requested answers reach the dashboard. It does not evaluate open-ended insight discovery or its broader analytical capabilities.

Vega-Lite defines concise visual encodings and interactions that compile into lower-level specifications \cite{vegalite}. Draco expresses visualization design knowledge as hard and soft constraints \cite{draco}. Our fixed contracts apply established ideas about constrained representations, not a new visualization grammar or solver. The issue is whether the model selects the intended analytical answer before a valid display is constructed.

Evaluation must also extend beyond syntactic success. VisEval distinguishes executable validity, satisfaction of chart and data requirements, and readability; its annotation design allows multiple acceptable charts rather than exact matching to one representation \cite{viseval}. This is particularly relevant to our unresolved scoring issue: correct requested means can appear in a supported claim type that a frozen predicate does not recognize. Our contribution is not the general observation that visualization quality is multidimensional. It is an artifact-grounded examination of numerical validity, longitudinal answer coverage, and compiler attribution in a tightly controlled task setting. Unlike VisEval's broader benchmark, the present comparison is small, developmental, and has no completed independent human adjudication.

W3C PROV-DM describes entities, activities, agents, and derivations \cite{prov}. Our distinction among evidence, selections, compiler additions, and later rendering applies these principles without claiming PROV conformance. It enables separate accounting for selected answers and supplied context. The contribution concerns responsibility and completeness, not the introduction of tool use, structured output, or provenance.
